\documentclass{article}
\usepackage{amsmath,amssymb,amsthm}
\usepackage{algorithm}
\usepackage{algpseudocode}
\usepackage{hyperref}
\usepackage{enumitem}
\newtheorem{theorem}{Theorem}
\newtheorem{lemma}{Lemma}
\newtheorem{assumption}{Assumption}
\newtheorem{definition}{Definition}
\newcommand{\norm}[1]{\left\lVert#1\right\rVert}
\newcommand{\ip}[2]{\left\langle #1,#2\right\rangle}
\begin{document}

\section*{Algorithm Name}
\textbf{Frank–Wolfe (Conditional Gradient) Method} for minimizing a smooth convex function over a compact convex set.

\section*{Mathematical Formulation}
Let 
\[
\min_{x\in C} f(x),
\]
where $f:\mathbb{R}^d\to\mathbb{R}$ is convex and $L$‑smooth, and $C\subset\mathbb{R}^d$ is compact and convex.  
At iteration $k$:

\begin{enumerate}[label=\arabic*.]
    \item Compute the gradient $g_k:=\nabla f(x_k)$.
    \item Solve the \emph{linear minimization oracle} (LMO) 
    \[
    s_k = \arg\min_{s\in C}\;\ip{g_k}{s}.
    \tag{LMO}
    \]
    \item Choose a stepsize $\gamma_k\in[0,1]$ (either a prescribed schedule $\gamma_k=\frac{2}{k+2}$ or an exact line‑search).  
    \item Perform the convex combination update 
    \[
    x_{k+1}= (1-\gamma_k)x_k + \gamma_k s_k .
    \tag{FW‑Update}
    \]
\end{enumerate}

The algorithm maintains all iterates inside $C$ because $C$ is convex and $\gamma_k\in[0,1]$.

\section*{Pseudocode}
\begin{algorithm}[H]
\caption{Frank–Wolfe}
\begin{algorithmic}[1]
\Require Convex set $C$, smooth convex $f$, initial point $x_0\in C$, stepsize rule $\{\gamma_k\}$ 
\For{$k=0,1,2,\dots$}
    \State $g_k \gets \nabla f(x_k)$
    \State $s_k \gets \displaystyle\arg\min_{s\in C}\ip{g_k}{s}$ \Comment{LMO}
    \State Choose $\gamma_k\in[0,1]$ (e.g. $\gamma_k=\frac{2}{k+2}$)
    \State $x_{k+1}\gets (1-\gamma_k)x_k+\gamma_k s_k$
    \If{stopping criterion met}
        \State \textbf{break}
    \EndIf
\EndFor
\State \Return $x_{k+1}$
\end{algorithmic}
\end{algorithm}

\section*{Dry Run Example}
Minimize $f(w)=(w-3)^2$ over the interval $C=[0,5]$.
The gradient is $\nabla f(w)=2(w-3)$.  
We initialise $w_0=0$ and use a fixed stepsize $\gamma=0.1$.

\begin{center}
\begin{tabular}{c|c|c|c|c}
$k$ & $w_k$ & $\nabla f(w_k)$ & $s_k=\arg\min_{s\in[0,5]}\ip{\nabla f(w_k)}{s}$ & $w_{k+1}$ \\ \hline
0 & $0$ & $-6$ & $5$ (since $\ip{-6}{s}$ is smallest at $s=5$) & $0.9\cdot0+0.1\cdot5=0.5$ \\
1 & $0.5$ & $-5$ & $5$ & $0.9\cdot0.5+0.1\cdot5=0.95$ \\
2 & $0.95$ & $-4.1$ & $5$ & $0.9\cdot0.95+0.1\cdot5=1.355$ \\
3 & $1.355$ & $-3.29$ & $5$ & $0.9\cdot1.355+0.1\cdot5=1.7195$
\end{tabular}
\end{center}
Objective values:
\[
f(w_0)=9,\; f(w_1)=6.25,\; f(w_2)=4.2025,\; f(w_3)=3.014.
\]
The function value decreases, illustrating the algorithmic steps.

\newpage
\section*{Assumptions}
\begin{assumption}[Convexity and Smoothness]\label{ass:convexsmooth}
The objective $f$ is convex and $L$‑smooth on $\mathbb{R}^d$, i.e.
\[
f(y)\le f(x)+\ip{\nabla f(x)}{y-x}+\frac{L}{2}\norm{y-x}^2,\qquad\forall x,y\in\mathbb{R}^d.
\]
\end{assumption}
\begin{assumption}[Compact Feasible Set]\label{ass:compact}
$C\subset\mathbb{R}^d$ is non‑empty, convex and compact. Define its \emph{diameter}
\[
D:=\max_{x,y\in C}\norm{x-y}<\infty .
\]
\end{assumption}
\begin{assumption}[Linear Minimization Oracle]\label{ass:lmo}
At each iteration the LMO (\ref{LMO}) can be solved exactly.
\end{assumption}

\section*{Main Convergence Theorem}
\begin{theorem}[Standard Frank–Wolfe Rate]\label{thm:fw}
Under Assumptions \ref{ass:convexsmooth}–\ref{ass:lmo}, let $\{x_k\}$ be generated by the Frank–Wolfe method with the stepsize schedule
\[
\gamma_k=\frac{2}{k+2},\qquad k\ge0.
\]
Define $f^\star:=\min_{x\in C}f(x)$. Then for all $k\ge0$,
\[
f(x_k)-f^\star\;\le\; \frac{2LD^2}{k+2}.
\tag{FW‑Rate}
\]
\end{theorem}
A corollary for a constant stepsize $\gamma\in(0,1]$ yields a linear convergence bound in the presence of strong convexity, but the theorem above suffices for the basic smooth convex case.

\section*{Theoretical Properties}
\begin{enumerate}[label=\arabic*.]
    \item \textbf{Stability.} Since each iterate is a convex combination of points in $C$, all iterates stay in $C$; thus the algorithm never leaves the feasible region.
    \item \textbf{Hyperparameter Sensitivity.} The canonical schedule $\gamma_k=2/(k+2)$ is parameter‑free and yields the $O(1/k)$ rate. Using a constant $\gamma$ results in a slower $O(1)$ suboptimality unless additional curvature (e.g., strong convexity) is present.
    \item \textbf{Comparison to Projected Gradient Descent (PGD).} FW avoids costly Euclidean projections at the price of a suboptimal $O(1/k)$ rate (vs.\ $O(1/k)$ for PGD with the same smoothness constant but with a smaller constant factor $L/2$ rather than $L$). When $C$ has a cheap LMO (e.g., a simplex or a trace‑norm ball), FW is computationally advantageous.
\end{enumerate}

\section*{Discussion}
The bound \eqref{FW‑Rate} depends only on the smoothness constant $L$ and the diameter $D$ of the feasible set. No explicit dependence on the dimension $d$ appears, making the method attractive for high‑dimensional structured domains where projections are expensive but linear minimization is cheap.

\newpage
\section*{Preliminaries (Definitions and Lemmas)}
\begin{definition}[Curvature Constant]\label{def:curvature}
For a convex, $L$‑smooth $f$ on $C$, define
\[
C_f := \sup_{\substack{x,s\in C\\\gamma\in[0,1]}}
\frac{2}{\gamma^2}\Bigl(f\bigl(x+\gamma(s-x)\bigr)-f(x)-\gamma\ip{\nabla f(x)}{s-x}\Bigr).
\]
\end{definition}
\begin{lemma}\label{lem:curvature-bound}
If $f$ is $L$‑smooth on $\mathbb{R}^d$, then $C_f\le LD^2$.
\end{lemma}
\begin{proof}
For any $x,s\in C$ and $\gamma\in[0,1]$, let $y=x+\gamma(s-x)$. By $L$‑smoothness,
\[
f(y)\le f(x)+\ip{\nabla f(x)}{y-x}+\frac{L}{2}\norm{y-x}^2
      =f(x)+\gamma\ip{\nabla f(x)}{s-x}+\frac{L}{2}\gamma^2\norm{s-x}^2 .
\]
Rearranging gives
\[
f(y)-f(x)-\gamma\ip{\nabla f(x)}{s-x}\le\frac{L}{2}\gamma^2\norm{s-x}^2
    \le\frac{L}{2}\gamma^2 D^2 .
\]
Multiplying by $2/\gamma^2$ yields the claim $C_f\le LD^2$.\qedhere
\end{proof}

\section*{Main Proof (Step‑by‑Step Derivation)}
We prove Theorem \ref{thm:fw}.  
All non‑trivial algebraic manipulations are numbered for easy reference.

\noindent\textbf{Notation.} Let $x^\star\in\arg\min_{x\in C}f(x)$. Define the \emph{Frank–Wolfe gap}
\[
g_k := \ip{\nabla f(x_k)}{x_k-s_k}\ge0 .
\tag{1}
\]

\begin{proof}
\textbf{[Step 1]} (Descent Lemma)  
By $L$‑smoothness with $y = x_{k+1}=x_k+\gamma_k(s_k-x_k)$,
\[
f(x_{k+1})\le f(x_k)+\gamma_k\ip{\nabla f(x_k)}{s_k-x_k}
               +\frac{L}{2}\gamma_k^2\norm{s_k-x_k}^2 .
\tag{2}
\]

\textbf{[Step 2]} (Replace inner product using the gap)  
Since $g_k = \ip{\nabla f(x_k)}{x_k-s_k}= -\ip{\nabla f(x_k)}{s_k-x_k}$,
\[
\ip{\nabla f(x_k)}{s_k-x_k}= -g_k .
\tag{3}
\]

\textbf{[Step 3]} (Bound the distance term)  
Because $C$ has diameter $D$, $\norm{s_k-x_k}\le D$. Hence
\[
\frac{L}{2}\gamma_k^2\norm{s_k-x_k}^2\le \frac{L}{2}\gamma_k^2 D^2 .
\tag{4}
\]

\textbf{[Step 4]} (Combine (2)–(4))  
\[
f(x_{k+1})\le f(x_k)-\gamma_k g_k+\frac{L}{2}\gamma_k^2 D^2 .
\tag{5}
\]

\textbf{[Step 5]} (Relate gap to suboptimality)  
From convexity,
\[
f(x_k)-f^\star \le \ip{\nabla f(x_k)}{x_k-x^\star}
                \le \ip{\nabla f(x_k)}{x_k-s_k}
                = g_k .
\tag{6}
\]
The first inequality follows because $x^\star\in C$ and $s_k$ minimizes the linear form over $C$.

\textbf{[Step 6]} (Insert (6) into (5)}  
\[
f(x_{k+1})-f^\star \le f(x_k)-f^\star-\gamma_k\bigl(f(x_k)-f^\star\bigr)
                     +\frac{L}{2}\gamma_k^2 D^2 .
\tag{7}
\]

\textbf{[Step 7]} (Define the error sequence $h_k:=f(x_k)-f^\star$)  
Equation (7) becomes
\[
h_{k+1}\le (1-\gamma_k)h_k+\frac{L}{2}\gamma_k^2 D^2 .
\tag{8}
\]

\textbf{[Step 8]} (Plug the canonical stepsize $\gamma_k=2/(k+2)$)  
\[
1-\gamma_k = 1-\frac{2}{k+2}= \frac{k}{k+2},\qquad
\gamma_k^2 = \frac{4}{(k+2)^2}.
\]
Thus (8) yields
\[
h_{k+1}\le \frac{k}{k+2}\,h_k + \frac{L}{2}\frac{4}{(k+2)^2}D^2
        = \frac{k}{k+2}\,h_k + \frac{2LD^2}{(k+2)^2}.
\tag{9}
\]

\textbf{[Step 9]} (Induction on $k$)  
We claim
\[
h_k\le \frac{2LD^2}{k+2},\qquad\forall k\ge0 .
\tag{10}
\]
Base case $k=0$: by definition $h_0=f(x_0)-f^\star\le \frac{2LD^2}{2}=LD^2$. Since $f$ is $L$‑smooth and $x_0,s_0\in C$, $h_0\le LD^2$ holds, establishing the base case.

Assume (10) holds for some $k$. Then from (9),
\[
h_{k+1}\le \frac{k}{k+2}\cdot\frac{2LD^2}{k+2}
          +\frac{2LD^2}{(k+2)^2}
        = \frac{2LD^2}{(k+2)^2}\bigl(k+1\bigr)
        = \frac{2LD^2}{(k+3)} .
\]
Since $k+3 = (k+1)+2$, this is exactly the bound (10) with index $k+1$. Hence the induction is complete.

\textbf{[Step 10]} (Conclusion)  
Inequality (10) is precisely the claim of Theorem \ref{thm:fw}. \qedhere
\end{proof}

\section*{Rate Derivation (With Constants)}
The proof above yields the explicit constant $2LD^2$ in the numerator.  
If the curvature constant $C_f$ from Definition \ref{def:curvature} is used instead of $LD^2$, the bound becomes
\[
f(x_k)-f^\star \le \frac{2C_f}{k+2},
\]
which is often tighter because $C_f\le LD^2$ (Lemma \ref{lem:curvature-bound}).

\section*{Special Cases}
\begin{enumerate}[label=\alph*.]
    \item \textbf{Strongly Convex $f$.} If $f$ is $\mu$‑strongly convex, then the Frank–Wolfe gap satisfies $g_k\ge \mu \norm{x_k-x^\star}^2$, and a \emph{linear} convergence rate can be obtained for a constant stepsize $\gamma\in(0,1]$ (see e.g., Jaggi 2013).
    \item \textbf{Exact Line Search.} Choosing $\gamma_k=\arg\min_{\gamma\in[0,1]} f\bigl(x_k+\gamma(s_k-x_k)\bigr)$ yields the same $O(1/k)$ bound but with a potentially smaller constant, as the term $\frac{L}{2}\gamma_k^2 D^2$ in (5) is replaced by the exact minimal value.
    \item \textbf{Polytope Domains.} When $C$ is a polytope, each LMO returns a vertex, and the iterates are sparse convex combinations of at most $k+1$ vertices. This sparsity property is a distinctive advantage over projected methods.
\end{enumerate}

\end{document}